\documentclass[a4paper,12pt]{letter}
% POZOR: samo utf8. V starih datotekah je bil za njim se \usepackage[cp1250]{inputenc},
% kar je prevzelo prednost in lahko pokvari sumnike.
\usepackage[utf8]{inputenc}
\usepackage[T1]{fontenc}
\usepackage[english,slovene]{babel}
\usepackage[tikz]{bclogo}
\usepackage{qrcode,marvosym}
\usepackage{pgf,tikz,pgfplots}
\pgfplotsset{compat=1.14}
\usepackage{mathrsfs}
\usetikzlibrary{arrows}
\usepackage{graphicx}
\newcommand{\ds}{\displaystyle}
\usepackage{fancybox}
\usepackage{amsfonts}
\usepackage{amsmath}
\usepackage{wallpaper}
\usepackage{hyperref}
\usepackage{array}
\newcolumntype{M}[1]{>{\centering\arraybackslash}m{#1}}
\newcolumntype{N}{@{}m{0pt}@{}}
\usepackage{fancyhdr,enumerate}
\newcommand{\ora}{\overrightarrow}
\renewcommand{\headrulewidth}{3pt}
\renewcommand{\headrule}{\hbox to\headwidth{%
  \color{gray}\leaders\hrule height \headrulewidth\hfill}}
\renewcommand{\footrulewidth}{2pt}
\renewcommand{\footrule}{\hbox to\headwidth{%
  \color{brown}\leaders\hrule height \footrulewidth\hfill}}
\usepackage{gensymb}
\usepackage{amssymb}
  \textwidth 20 true cm
  \oddsidemargin -2 true cm
  \evensidemargin -2 true cm
  \topmargin -3.5 true cm
  \textheight 27.5 true cm
  \linespread{1.2}
\renewcommand{\headrulewidth}{0.3pt}
\renewcommand{\footrulewidth}{1.9pt}
\definecolor{qqwuqq}{rgb}{0.,0.39215686274509803,0.}
\definecolor{qqqqff}{rgb}{0,0,1}
\definecolor{ffqqqq}{rgb}{1.,0.,0.}
\definecolor{zzuxux}{rgb}{0,0,0}
\usepackage{mdframed}
\usepackage{multicol}
\definecolor{bgblue}{RGB}{0,0,253}
\usepackage{eurosym}
%%%%%%%%%%%%%%%%%%%%%%%%%%%%%%%%%%%%%%%%%%%%%%%%%%%%%%%%
\newcounter{tocke}
\setcounter{tocke}{0}
\usepackage{tikz-3dplot}
%%%%%%%%%%%%%%%%%%%%%%%%%%%%%%%%%%%%%%%%%%%%%%%%%%%%%%%%%%
\mdfdefinestyle{theoremstyle}{%
linecolor=brown!40,linewidth=1pt,%
frametitlerule=true,%
frametitlebackgroundcolor=brown!50,
innertopmargin=\topskip,
nobreak=true,
}
\mdfdefinestyle{theoremstyle1}{%
linecolor=brown,middlelinewidth=1pt,%
frametitlerule=true,%
apptotikzsetting={\tikzset{mdfframetitlebackground/.append style={%
shade,left color=black!40, right color=gray!10},
mdfbackground/.append style={
top color=white!100!white,
bottom color=black!5
},
}},
frametitlerulecolor=gray,
frametitlerulewidth=2pt,
innertopmargin=\topskip,
innerbottommargin=\topskip,
nobreak=true,
}
\mdtheorem[style=theoremstyle1]{naloga}{Naloga}
\mdtheorem[style=theoremstyle]{resitev}{Re\v{s}itev}
\mdtheorem[style=theoremstyle1]{kriterij}{Kriterij ocenjevanja}
%%%%%%%%%%%%%%%%%%%%%%%%%%%%%%%%%%%%%%%%%%%%%%%%%%%%%%%%%%
\usepackage[table]{xcolor}
\usepackage[printwatermark]{xwatermark}
\newwatermark[allpages,color=brown!2,angle=45,scale=5,
xpos=-5,ypos=0]{matej.info}
%%%%%%%%%%%%%%%%%%%%%%%%%%%%%%%%%%%%%%%%%%%%%%%%%%%%%%%%%%

\begin{document}
\pagestyle{fancy}
\renewcommand\headrulewidth{0pt}
\lhead{}\chead{}\rhead{}
\rfoot{\vspace*{-2\baselineskip}}
\renewcommand{\vec}{\overrightarrow}
\everymath{\displaystyle}

\begin{minipage}{20cm}
\begin{bclogo}[couleur=brown!40,
  arrondi=0,ombre=true, couleurOmbre=gray!50,
logo =\bcplume,
  noborder=true]
{\textrm{{\Large{TEST 2.N.2} \Huge{}:
\large{(algebrski) Ulomki, decimalna števila, procent, sorazmerja, enačbe, neenačbe} \hfill $PTI-1$}}}
\end{bclogo}
\end{minipage}
\begin{minipage}{20cm}
\begin{bclogo}[couleur=brown!40,ombre=true, couleurOmbre=gray!50,
  arrondi=0,
logo =\bcplume,
  noborder=true]
{\textsc{{\large{Ime in priimek:}
\vspace{0.2cm} \rule{15cm}{0.2mm}\hfill }}}
\end{bclogo}
\end{minipage}




%%%%%%%%%%%%%%%%%%%%%%%%%%%%%%%%%%%%%%%%%%%%%%%%%
\begin{naloga}[\hfill \textit{ve\v{c} operacij v istem ra\v{c}unu}]
Izra\v{c}unaj. Upo\v{s}tevaj vrstni red ra\v{c}unskih operacij.

\begin{multicols}{2}
\begin{enumerate}[a.)]
\item $\dfrac{3}{4}+\dfrac{1}{2}\cdot\dfrac{2}{3}$
\item $\left(\dfrac{2}{3}-\dfrac{1}{4}\right):\dfrac{5}{6}$
\end{enumerate}
\end{multicols}
\end{naloga}

%%%%%%%%%%%%%%%%%%%%%%%%%%%%%%%%%%%%%%%%%%%%%%%%%
\begin{naloga}[\hfill \textit{kraj\v{s}anje algebrskih ulomkov}]
Okraj\v{s}aj algebrski ulomek. \v{S}tevec in imenovalec najprej razstavi.

\begin{multicols}{2}
\begin{enumerate}[a.)]
\item $\dfrac{x^2-25}{x-5}$
\item $\dfrac{3x^2-12}{x+2}$
\item $\dfrac{x^2+x-6}{x^2-4}$
\item $\dfrac{2x^2-6x}{x^2-9}$
\end{enumerate}
\end{multicols}
\end{naloga}

%%%%%%%%%%%%%%%%%%%%%%%%%%%%%%%%%%%%%%%%%%%%%%%%%
\begin{naloga}[\hfill \textit{ena\v{c}be z ulomki}]
Re\v{s}i ena\v{c}bo.
$$\dfrac{2x-1}{3}=\dfrac{x+3}{2}$$
\end{naloga}

%%%%%%%%%%%%%%%%%%%%%%%%%%%%%%%%%%%%%%%%%%%%%%%%%
\begin{naloga}[\hfill \textit{besedilne naloge z ulomki}]
Zapi\v{s}i ra\v{c}un oziroma ena\v{c}bo in re\v{s}i.

\begin{enumerate}[1.)]
\item Marko je prvi dan prekolesaril $\dfrac{3}{8}$ poti, drugi dan pa \v{s}e $\dfrac{1}{4}$ poti.
Kolik\v{s}en del poti je prekolesaril v teh dveh dneh? Kolik\v{s}en del mu je \v{s}e ostal?
\item Babica je $\dfrac{2}{5}$ svojih prihrankov podarila vnuku, preostanek pa vnukinji.
Prihranki zna\v{s}ajo $3500$~\euro{}. Koliko \euro{} je prejela vnukinja?
\item Steklenica drži $\dfrac{3}{4}$ litra. Koliko steklenic napolnimo z $9$ litri soka?
\end{enumerate}
\end{naloga}

%%%%%%%%%%%%%%%%%%%%%%%%%%%%%%%%%%%%%%%%%%%%%%%%%
\begin{naloga}[\hfill \textit{premo sorazmerje}]
Re\v{s}i nalogo s premim ali obratnim sorazmerjem.

Stroj napolni $840$ steklenic v $35$ minutah. Koliko minut potrebuje za $1320$ steklenic?
\end{naloga}

%%%%%%%%%%%%%%%%%%%%%%%%%%%%%%%%%%%%%%%%%%%%%%%%%
\begin{naloga}[\hfill \textit{besedilne naloge -- odstotki}]
Zapi\v{s}i ra\v{c}un oziroma ena\v{c}bo in re\v{s}i.

Cena kolesa se je zvi\v{s}ala za $12\%$ in zdaj zna\v{s}a $504$~\euro{}. Kolik\v{s}na je bila cena kolesa pred podra\v{z}itvijo?
\end{naloga}

%%%%%%%%%%%%%%%%%%%%%%%%%%%%%%%%%%%%%%%%%%%%%%%%%
\begin{naloga}[\hfill \textit{periodi\v{c}na \v{s}tevila}]
Pretvori med ulomkom in decimalnim (periodi\v{c}nim) zapisom \v{s}tevila.

\begin{multicols}{2}
\begin{enumerate}[a.)]
\item Zapi\v{s}i $0.\overline{7}$ kot okraj\v{s}an ulomek.
\item Zapi\v{s}i $0.\overline{36}$ kot okraj\v{s}an ulomek.
\item Zapi\v{s}i $\dfrac{7}{12}$ kot decimalno \v{s}tevilo.
\item Zapi\v{s}i $0.2\overline{3}$ kot okraj\v{s}an ulomek.
\item Izra\v{c}unaj $0.\overline{2}+0.\overline{5}$ tako, da \v{s}tevili najprej pretvori\v{s} v ulomka.
\end{enumerate}
\end{multicols}
\end{naloga}

%%%%%%%%%%%%%%%%%%%%%%%%%%%%%%%%%%%%%%%%%%%%%%%%%
\begin{naloga}[\hfill \textit{me\v{s}ano -- periodi\v{c}na \v{s}tevila in potence s celimi eksponenti}]
Izra\v{c}unaj. Periodi\v{c}na \v{s}tevila najprej pretvori v ulomke.

\begin{multicols}{2}
\begin{enumerate}[a.)]
\item $\left(2^{-1}+0.\overline{7}\right):\dfrac{1}{18}$
\item $\left(0.\overline{6}+4^{-1}\right):\dfrac{11}{24}$
\end{enumerate}
\end{multicols}
\end{naloga}

\vfill
\begin{center}
\qrcode[hyperlink,height=0.8in]{http://www.matej.info}
\end{center}

\newpage
% Razred "letter" ne pozna ukaza \section, zato naslov napisemo neposredno.
\begin{center}
{\Large\textbf{Kon\v{c}ne re\v{s}itve}}
\end{center}
\vspace{0.3cm}

\begin{enumerate}

\item a.) $\dfrac{3}{4}+\dfrac{1}{3}=\dfrac{13}{12}$ \quad b.) $\dfrac{5}{12}\cdot\dfrac{6}{5}=\dfrac{1}{2}$

\item a.) $x+5$ \quad b.) $3x-6$ \quad c.) $\dfrac{x+3}{x+2}$ \quad d.) $\dfrac{2x}{x+3}$

\item $4x-2=3x+9 \Rightarrow x=11$

\item 1.) prekolesaril $\dfrac{5}{8}$ poti, ostane \v{s}e $\dfrac{3}{8}$ \quad
2.) $\dfrac{3}{5}\cdot3500=2100$~\euro{} \quad
3.) $9:\dfrac{3}{4}=12$ steklenic

\item $55$ min

\item $504:1.12=450$~\euro{}

\item a.) $0.\overline{7}=\dfrac{7}{9}$ \quad b.) $0.\overline{36}=\dfrac{4}{11}$ \quad
c.) $\dfrac{7}{12}=0.58\overline{3}$ \quad d.) $0.2\overline{3}=\dfrac{7}{30}$ \quad
e.) $\dfrac{2}{9}+\dfrac{5}{9}=\dfrac{7}{9}$

\item a.) $\left(\dfrac{1}{2}+\dfrac{7}{9}\right):\dfrac{1}{18}=\dfrac{23}{18}\cdot18=23$ \quad
b.) $\left(\dfrac{2}{3}+\dfrac{1}{4}\right):\dfrac{11}{24}=\dfrac{11}{12}\cdot\dfrac{24}{11}=2$

\end{enumerate}
\end{document}
